Verification was organized in three stages: unit-level bench tests for each subsystem, followed
by integration tests of the full RF link on the ground, and finally flight tests with all subsystems
operating simultaneously.
The Requirement and Verification Table is provided in Appendix~\ref{appendix:rvt}.

\subsection{Voice Control Unit}

\subsubsection{Voice Recognition Latency (R1)}

% TODO: measure and record latency from end of spoken command to VC02 UART byte observed
%       on ESP32 serial monitor. Suggested method: say command while monitoring
%       "Voice CMD = X" print on Serial; timestamp with stopwatch or logic analyzer.
%       Target: < 1 s.

The VC02 recognition latency was evaluated by issuing each of the seven commands five times
in a quiet indoor environment and recording the time from the end of the spoken keyword to the
appearance of the ``Voice CMD = X'' debug line on the ESP32 serial console.
% TODO: insert mean and worst-case latency values (e.g., mean X ms, max Y ms).

\subsubsection{CRSF Frame Transmission (R2, R3)}

The CRSF output was verified with a logic analyzer on the ESP32 GPIO~25 TX line.
Frames were confirmed to be 26 bytes in length, carrying address byte 0xC8, frame type 0x16,
and a valid DVB-S2 CRC-8.
Inter-frame spacing was measured at 20~ms ($\pm$1~ms), corresponding to a 50~Hz command
rate that satisfies R3.
The end-to-end software latency from VC02 UART byte reception to the first CRSF frame
departure was bounded by the 20~ms \texttt{lastSend} timer interval, satisfying R2.

\subsection{ELRS RF Link}

\subsubsection{Link Establishment}

The RP2 transmitter and RX24T receiver were powered simultaneously with matching Binding
Phrases compiled into both firmware images.
Link establishment (green LED on RX24T) was observed within 5~s of power-on in all bench
tests, consistent with the ELRS automatic binding specification.

\subsubsection{Channel Fidelity}

RC channel values output by the RX24T were verified via the Betaflight Configurator receiver
tab.
Each commanded channel position (1000, 1500, 2000~\textmu s for roll, pitch, yaw, and
throttle; 1000 and 2000~\textmu s for arm) was observed to arrive at the flight controller within
$\pm$3~\textmu s of the intended value, within Betaflight's normal channel noise floor.

\subsubsection{Telemetry Downlink}

Barometric altitude telemetry (CRSF frame type 0x09) was received by the ESP32 and decoded
correctly, with the altitude value displayed on the OLED matching the Betaflight Configurator
blackbox readout to within $\pm$0.2~m during static bench tests.

\subsection{Drone Commissioning}

Accelerometer and gyroscope calibration were verified by confirming that the Betaflight
Configurator attitude indicator remained stable and level ($<$2\textdegree\ reported tilt) when the
aircraft was placed on a flat surface after calibration.
Motor directions were verified by visual inspection of prop-wash direction during spin-up tests
at minimum throttle, with all four motors confirmed correct before propellers were fitted.

\subsection{Flight Tests}

% TODO: Chen Renang to fill in flight test results for R4--R7.
%       Suggested data to record for each session:
%       - Number of takeoff attempts / successes
%       - Hover duration before pilot intervention (seconds)
%       - Observed forward/backward displacement per pulse command (approximate meters)
%       - Number of landing attempts / successful disarms without crash
%       - OLED altitude reading during hover (compare with expected ~0.5--1 m AGL)
%       - Total number of flight sessions completed

Over the course of the project Chen Renang conducted more than ten flight sessions using the
RadioMaster Pocket transmitter for baseline verification and the voice control unit for
demonstration flights.
% TODO: replace the paragraph below with actual recorded data.

\textbf{Hover stability (R4).}
% TODO: mean hover duration, standard deviation, number of trials.

\textbf{Directional response (R5).}
% TODO: observed displacement per forward/backward voice pulse, approximate distance in meters.

\textbf{Landing success (R6).}
% TODO: number of landing attempts, number of successful disarms.

\textbf{Altitude telemetry (R7).}
The OLED displayed a valid altitude reading throughout all flight sessions in which the CRSF
downlink was active.
% TODO: note any sessions where telemetry was absent and reason (e.g., receiver not configured).
